% !TeX root = ../main.tex

\xchapter{任务概述与相关理论}{Task Overview and Related Theory}

本章围绕平面几何问题求解任务，对本文后续研究所涉及的基本问题和相关技术进行说明。首先，从任务层面对平面几何问题求解的类型、形式化表示和特点进行介绍，并对本文使用的“长链条”概念作出统一定义；其次，从技术层面对后续研究中涉及的大语言模型、微调方法和强化学习方法进行概述，重点说明这些方法的基本原理及其在复杂推理任务中的一般作用。

\xsection{任务定义}{Task Definition}

本文关注平面几何问题求解任务（Plane Geometry Problem Solving, PGPS）。该任务以题目文本、几何图像和待求目标为输入，要求模型结合题设条件与图形结构，生成相应的推理过程并输出最终答案。本文的研究范围限定在中学课程标准下的欧氏平面几何，不涉及解析几何、立体几何等其他几何分支。如图~\ref{fig:task_difference}所示，与侧重坐标表示与代数运算的解析几何、强调空间关系构建的立体几何不同，平面几何问题求解主要围绕二维图形中的点、线、角、圆等基本对象展开，更关注几何对象之间的关系识别、条件约束理解以及基于定义、公理和定理的逻辑推导过程。

\begin{figure}[htbp]
    \centering
    \includegraphics[width=1.0\textwidth]{c2_difference.png}
    \caption{平面几何、解析几何与立体几何题目示例}
    \label{fig:task_difference}
\end{figure}

为了对该任务进行统一描述，在此从形式化角度对其输入、输出和求解目标进行建模。设题目文本序列为 $T$，原始几何图像为 $D \in \mathbb{R}^{H \times W \times C}$，待求目标为 $Q$，则几何问题求解系统需要根据输入生成推理序列 $S$ 并输出最终答案 $A$。记由参数 $\theta$ 构建的求解模型为 $M_{\theta}$，其任务是学习题目条件、图形结构与求解目标之间的映射关系。该过程可统一表示为如下条件生成问题：
\begin{equation}
    (S^*, A^*) = \operatorname*{arg\,max}_{S \in \mathcal{S}, A \in \mathcal{A}} P(S, A \mid T, D, Q; \theta),
\end{equation}
其中，$\mathcal{S}$ 表示合法的欧氏平面几何推理序列空间，$\mathcal{A}$ 表示候选答案空间。

上述定义表明，平面几何问题求解不仅要求模型输出最终结果，还要求其在图文联合条件下组织中间推理过程。为进一步聚焦本文所研究的复杂推理场景，下面首先对长链条问题进行界定，并在此基础上总结该任务的主要特点。

\xsubsection{长链条问题定义}{Definition of Long-chain Problems}

本文将长链条几何问题定义为：在人类专家的标准解题过程中，需要至少 $5$ 个逻辑推理步骤才能完成求解的问题。记标准推理序列为 $S=\{s_1,s_2,\dots,s_n\}$，其中 $n$ 为逻辑推理步骤数，则当 $n \geq 5$ 时，将其视为长链条几何问题。

这里的“逻辑推理步骤”是指依据已有条件或前一步结论，结合几何定义、性质、公理或定理推出新结论的过程；单纯的数值运算或未产生新结论的描述不单独计步。后文将以逻辑推理步骤数作为衡量题目推理深度的主要指标。

\xsubsection{任务特点}{Task Characteristics}

平面几何问题通常以题目文本与辅助图形共同呈现，已知条件和几何关系分布在文字描述与图像结构之中。尽管也存在纯文本几何任务，但在中学教学与考试场景下，图文结合仍是主流形式，因此本文将其视为典型的图文联合多模态推理任务。从题目呈现形式来看，平面几何问题常见于选择题、填空题和解答题。不同题型在答案表达的完整性上存在差异，但并不改变任务本身的核心性质。无论题目以何种形式出现，模型都需要结合题目条件与图形结构，围绕求解目标组织推理过程。

按照求解目标，平面几何问题可分为数值计算类和结论证明类。前者关注长度、角度、面积等结果求解，后者关注几何关系的判断与证明。两类问题虽然输出形式不同，但都依赖于对题设条件、图形结构和中间结论的综合理解。此外，平面几何问题的求解通常不是对最终答案的直接预测，而需要经过一定的中间推导过程，因此任务评价除关注结果正确性外，也应兼顾推理过程的合理性与一致性。

\xsection{相关技术介绍}{Related Technology Introduction}

本节介绍与平面几何问题求解相关的几类基础技术，主要包括大语言模型、微调方法和强化学习方法。这些技术已经在自然语言处理、多模态理解和复杂推理等任务中得到广泛研究。

\xsubsection{大语言模型}{Large Language Models}

大语言模型（Large Language Models, LLMs）是在预训练语言模型基础上发展起来的一类模型，其技术起点通常追溯到 Transformer ~\cite{vaswani2017attention}架构的提出。相比早期依赖循环神经网络和长短时记忆网络的序列建模方法，Transformer 采用自注意力机制，可以直接建模序列中不同位置之间的关系，并支持并行训练，因此很快成为自然语言处理中的主流结构。在这一基础上，BERT~\cite{devlin2019bert}、GPT~\cite{radford2018gpt}等模型推动了预训练语言模型的发展。BERT 采用双向编码方式，更适合理解类任务；GPT 采用自回归生成方式，更适合文本生成任务。随着参数规模、训练数据和训练策略不断扩展，语言模型逐步发展为具备较强知识表示、上下文理解和文本生成能力的大语言模型。

从基本原理上看，大语言模型的核心仍然是语言建模，即在大规模语料上学习句子和篇章中的统计规律与语义关系。给定输入后，模型先将离散词元映射为连续表示，再通过多层自注意力和前馈网络不断更新上下文表示，最后完成下一个词元预测或被遮蔽词元恢复。自注意力机制通过查询、键和值之间的匹配关系分配权重，使模型能够根据上下文动态关注不同位置的信息；多头注意力则从不同子空间同时建模语义关系，从而增强表示能力。由于这种统一的建模方式，大语言模型能够迁移到问答、摘要和推理等多种任务中，已经成为当前自然语言处理的重要基础模型~\cite{brown2020language}。

除纯文本模型外，多模态大语言模型（Multimodal Large Language Models, MLLMs）进一步将图像等非文本模态纳入统一框架。这类模型通常先利用视觉编码器提取图像特征，再通过投影层或跨模态对齐模块将视觉表示映射到语言空间，最后由语言模型统一完成理解与生成~\cite{liu2024improved}。这一技术路线并没有改变语言模型作为核心的地位，而是在输入端增加了视觉信息接口。随着模型能力不断提升，多模态大语言模型已经可以较好地处理图文问答、图像描述和视觉推理等任务。如图 \ref {fig:llm_survey} 所示，近年来一批能力强劲的大语言模型相继涌现，它们普遍支持多模态输入，可处理图像、文本、文档、音频等多种类型的信息。

近年来，大语言模型在复杂推理任务上的能力也受到越来越多关注，进而出现了“思考型”（Thinking）大语言模型的研究。这类方法的重点不是改变底层结构，而是在训练或推理阶段显式引入中间推导过程，使模型先给出推理路径，再生成最终答案。这一思路与链式思维提示（Chain-of-Thought, CoT）~\cite{wei2022cot}密切相关。已有研究表明，当模型被引导输出中间推理过程时，其在数学推理、常识推理和符号推理等任务上的表现通常会有所提升。因此，大语言模型的发展已不再局限于一般文本生成，而是逐步扩展到对复杂推理能力的建模与增强。

\begin{figure}[htbp]
    \centering
    \includegraphics[width=1.0\textwidth]{c2_LLM_survey.png}
    \caption{大语言模型发展时间线~\cite{zhao2023llmsurvey}}
    \label{fig:llm_survey}
\end{figure}

\xsubsection{微调方法}{Fine-tuning Methods}

预训练语言模型虽然具有较强的通用能力，但并不是直接针对某一具体任务设计的。为了使模型适应新的任务形式和输出要求，通常需要在预训练模型基础上继续使用目标任务数据进行训练，这一过程称为微调（Fine-tuning）。随着预训练语言模型的发展，自然语言处理中的训练范式也逐渐从“针对单一任务从头训练”转向“先预训练、后微调”。这种方式减少了对大规模标注数据的依赖，也提高了模型在下游任务中的适应能力~\cite{brown2020language}。

指令监督微调（Instruction Supervised Fine-tuning）是在一般微调方法基础上的进一步发展。它将任务统一写成“指令—输入—输出”的形式，使模型学习按照自然语言指令完成不同任务。与传统分类或序列标注方式相比，指令微调更接近生成式模型的使用方式，也更适合当前的大语言模型。其基本做法是将任务说明、上下文和目标答案组织成训练样本，并通过监督学习优化模型输出与参考答案之间的差异。研究表明，经过指令微调的模型通常具有更好的指令遵循能力，在零样本和少样本场景下也表现出较强的泛化能力~\cite{wang2023selfinstruct}。


随着模型规模不断增大，全参数微调的代价越来越高。它不仅需要较大的显存和存储资源，也会增加训练和部署成本。为此，研究中提出了参数高效微调（Parameter-Efficient Fine-tuning, PEFT）方法，其基本思想是在尽量冻结原始模型参数的前提下，只训练少量新增参数来完成任务适配。LoRA~\cite{hu2022lora}（Low-Rank Adaptation）是其中最有代表性的方法之一，其核心做法是在原有线性变换矩阵上加入低秩更新项，用两个较小矩阵的乘积来近似参数变化，而不是直接更新整个权重矩阵。这样可以在保持原模型主体参数不变的情况下，以较少的可训练参数完成适配。

从原理上看，LoRA 假设任务适配所需的参数变化主要集中在一个低维子空间内，因此无需对全部参数进行独立更新，而可以通过少量新增参数实现模型适配。与全参数微调相比，LoRA 所需训练参数更少，显存开销更低，也更便于保存和切换不同任务的适配模块，因此在保证一定适配效果的同时显著降低了训练与部署成本。由于兼具效率和灵活性，LoRA 已成为当前大语言模型微调中的常用方法。

\xsubsection{强化学习}{Reinforcement Learning}

强化学习（Reinforcement Learning, RL）是一类通过智能体与环境交互来学习策略的方法。与监督学习直接依赖固定标签不同，强化学习中的模型需要根据当前状态选择动作，环境再返回奖励和新的状态，模型据此不断更新策略，以最大化长期累计回报~\cite{sutton1998reinforcement}。从形式上看，强化学习通常可表示为一个马尔可夫决策过程（Markov Decision Process, MDP），其核心要素包括状态空间、动作空间、状态转移概率和奖励函数。对于给定策略 $\pi$，目标是最大化期望回报
\begin{equation}
    J(\pi)=\mathbb{E}_{\tau \sim \pi}\left[\sum_{t=0}^{T}\gamma^t r_t\right],
\end{equation}
其中 $\tau$ 表示由策略生成的轨迹，$\gamma$ 为折扣因子，$r_t$ 表示第 $t$ 步的奖励。

对于生成式模型而言，强化学习可以将文本生成过程看作一个序列决策过程。给定输入提示后，模型每生成一个词元或一个推理片段，都可以看作在当前状态下执行一次动作；整个输出序列则构成一条轨迹。这样一来，模型优化的目标就不再只是拟合参考答案，而是能够直接根据整条输出的质量进行更新。因此，强化学习特别适用于那些难以通过单一监督标签刻画、却可以通过整体反馈进行评价的任务，例如回答质量、推理正确性和偏好对齐等。

在大语言模型的强化训练中，常见做法是采用策略梯度方法对生成策略进行更新。其中，PPO（Proximal Policy Optimization）是较有代表性的策略优化方法。它的基本思想是在更新策略时限制新旧策略之间的变化幅度，避免参数更新过大导致训练不稳定。该方法通过对策略更新施加局部约束，在优化奖励的同时控制策略分布的过快漂移，因此在大模型对齐和推理优化中得到了广泛应用~\cite{schulman2017ppo,ouyang2022training}。

在此基础上，面向大模型推理训练又进一步发展出了组相对策略优化（Group Relative Policy Optimization, GRPO）方法~\cite{guo2025deepseekr1}。与传统方法主要依赖绝对奖励和显式价值估计不同，GRPO 更关注同一输入下多条候选输出之间的相对优劣关系。具体来说，给定一个问题，模型会采样得到一组候选轨迹，再根据组内比较结果构造相对优势，并据此更新策略。

从原理上看，GRPO 的关键不在于单条输出获得了多少绝对分数，而在于它在同组候选中的相对表现是否更好。这样做可以减弱不同样本之间奖励尺度不一致带来的影响，也能够降低训练过程对单独价值模型的依赖。与传统策略优化方法相比，GRPO 更适合用于推理任务，因为在这类任务中，不同候选解答之间的优劣通常更容易比较，而不容易获得稳定的绝对奖励信号。与此同时，GRPO 在训练过程中通常仍会引入参考策略来约束当前策略不要偏离初始模型过远，这种约束一般可以通过 KL 散度来刻画。KL 散度的作用在于控制策略漂移，使模型在追求更高奖励的同时保持输出分布的基本稳定性。

总体来看，GRPO 可以理解为一种结合组内相对比较与策略分布约束的优化方法。前者提供了更适合推理任务的训练信号，后者则保证训练过程的稳定性。对于大语言模型尤其是推理模型来说，GRPO 能够较好地利用多条候选解答之间的相对差异，从而为复杂推理任务提供更有效的强化学习训练机制。

\xsection{本章小结}{Chapter Summary}

本章对平面几何问题求解任务及其相关技术基础进行了概述。首先，围绕任务本身，介绍了平面几何问题的基本类型、任务特点、长链条问题定义及其形式化表示，为后续数据构建与方法设计提供了统一的问题描述。其次，围绕相关技术，介绍了大语言模型、微调方法和强化学习的基本原理，并进一步说明了多模态大语言模型、指令监督微调、LoRA 微调以及 GRPO 等方法的核心思想。以上内容构成了本文后续章节展开的理论基础。 